\chapter{What a virtual laboratory can establish}
\label{vl:entry}

A proof identifies what follows from an experimental contract. A virtual
laboratory asks whether a proposed apparatus and analysis actually satisfy
that contract, whether the signal survives realistic forms of imperfection,
and what a failed experiment would look like. It cannot establish that an
unmeasured physical apparatus follows the simulated equations.

This extension supplies three connected laboratories. A driven-ring experiment
generates individual physical jump counts and noisy endpoint readings. A qubit
experiment propagates prepared states through independently implemented Bloch
and Lindblad dynamics, with preparation, control and detector errors. A
metastable classical experiment generates a chronological series of dependent
pairs, exposing the difference between stationary marginal probabilities and
independent records. The inference routines receive their declared counts and
calibrations, not the hidden generator used to score performance.

The simulations have already changed the theory and protocol design. Timing
noise sometimes preserves the exact comparison structure, while an apparently
innocent exposure convention can create false equilibrium violations. A
conditional-mean envelope corrects dependent sampling without an unjustified
effective sample size. A fixed polynomial contrast uses the sampling budget
more effectively for a known target, while the general simultaneous-box
procedure retains greater freedom to choose targets after acquisition.
A common clock correction repairs phase mismatch in the dimension experiment;
its calibration expense must be compared with simply collecting more records.

\section{The evidence contract}
Each study declares its data-generating model, target, acquisition procedure,
inference rule, cost accounting and simulation repetitions. Frozen JSON files
and seeded streams make the reported samples reproducible. Pilot and held-out
experiments use separate streams. Refinement loops are identified explicitly:
a revised protocol is frozen before its new evaluation data are generated.
The simulations retain unfavorable regimes and competing explanations.

The reported error guarantees and Monte Carlo intervals answer different
questions. A theorem may guarantee at most $5\%$ erroneous inference under its
premises. Repeated simulations estimate the actual operating probability in a
specified model, with finite Monte Carlo uncertainty. Zero errors among a
finite number of trials does not establish zero error probability. A numerical
solver agreement checks implementation, while a mathematical confidence bound
checks an inference rule; neither validates a real detector.

All time units and physical rates in this extension are illustrative. Entropy
rates are reported in natural units per model time. Converting them into watts
requires the thermodynamic embedding, temperature and time calibration specified
earlier. Evolution time, pilot measurements, reset and detector overhead are
reported separately. Unmodeled overhead is not assigned zero physical cost.

\section{How to use the extension}
Begin with the timing chapter before applying a correlation test to finite
exposures. Use the sampling chapter when records share a physical trajectory.
For a predeclared target, the contrast chapter provides a direct finite-record
certificate; for later exploratory selection, retain the simultaneous-box
workflow. The two apparatus chapters report complete simulation designs and
controls. The final comparison chapter separates practical gains from failures
and connects them to the next independent physical validation.

The software package includes exact certificates, numerical simulators,
retained records, analysis code, independent checks and reproduction commands.
It is intended to let another scientist replace a model or apparatus contract,
rerun the same falsification tests and determine whether a claimed conclusion
survives that change.
